\section{Kết luận}\label{sec:conclusion}

Báo cáo đã tìm hiểu hai công trình tiếp cận vấn đề phối hợp đa tác tử từ hai phía. Agent-Radar quản lý \emph{thông tin} trong giao tiếp: thay vì nén hoặc cắt lịch sử trao đổi, nó chấm điểm từng câu theo độ liên quan ngữ nghĩa nhân với suy giảm theo khoảng cách trên đồ thị và theo thời gian, rồi lái attention của LLM về phần được chọn. DAIN quản lý \emph{thành phần} tham gia suy luận: một Meta-Controller kích hoạt thưa nhóm tác tử chuyên biệt theo từng mẫu đa phương thức và điều phối trao đổi nén giữa chúng, huấn luyện end-to-end với loss đa mục tiêu. Điểm gặp nhau của hai thiết kế là lựa chọn phụ thuộc ngữ cảnh và sử dụng cấu trúc liên kết như một tín hiệu về độ tin cậy của thông tin; điểm khác cơ bản là một bên là lớp phủ training-free lúc suy luận, một bên là kiến trúc học từ đầu.

Phần thực nghiệm kiểm chứng được những gì tài nguyên cho phép. Với Agent-Radar, các công thức chọn ngữ cảnh được hiện thực đúng và đo được trên HotpotQA: bộ chọn giữ lại bằng chứng với mật độ cao hơn hẳn so với dùng cả lịch sử, với số token nhỏ hơn hàng chục lần; đồng thời thí nghiệm cũng cho thấy khi cơ chế chọn được dùng như bộ lọc cứng (không có bước lái attention trên toàn bộ lịch sử như paper gốc), chất lượng trả lời của một bộ đọc yếu giảm xuống -- kết quả này trái chiều với trục chính của paper nhưng nhất quán với lập luận của chính paper rằng chọn ngữ cảnh phải đi kèm cách đưa nó vào lúc suy luận. Với DAIN, bản đơn giản hóa huấn luyện được trên dữ liệu thật ENRICO và tái hiện đúng thứ tự ablation trong paper (mô hình đầy đủ tốt hơn biến thể bỏ giao tiếp, rồi đến biến thể bỏ điều phối); tuy nhiên với dữ liệu nhỏ và encoder nhẹ, mức tăng của DAIN so với early fusion trong paper chưa xuất hiện trong thiết lập của chúng tôi.

Những giới hạn còn tồn tại cần nói rõ. Thí nghiệm 1 là mô phỏng lịch sử trao đổi bằng quy tắc, không phải hành vi của LLM; môi trường mô phỏng cũng không tái tạo được tương quan giữa vị trí không gian/thời gian và độ tin cậy của thông tin -- chính yếu tố mà hai hệ số suy giảm của Agent-Radar dựa vào để phát huy tác dụng. Thí nghiệm 2 thay encoder lớn bằng encoder nhẹ và chỉ chạy một trong năm bộ dữ liệu, nên số liệu không so sánh trực tiếp với paper. Cả hai paper đều chưa công bố mã nguồn để đối chiếu từng dòng. Hướng phát triển tự nhiên tiếp theo là chạy phần lái attention (SPA) trên một LLM cục bộ nhỏ để khép lại vòng đánh giá cho Agent-Radar, và mở rộng thí nghiệm DAIN sang bộ dữ liệu thứ hai với encoder mạnh hơn.
